% theme: information_communication_network
\MajorQuestion{情報通信ネットワーク}
\SetRowSpaceQanda

\Instruction{ネットワークの基本に関する次の問いに答えなさい。}
\QQ{情報機器どうしを有線や無線でつなぎ、情報のやりとりができるようにしたものを何というか。}{情報通信ネットワーク}
\QQ{世界規模の通信ができる、代表的な情報通信ネットワークを何というか。}{インターネット}
\QQ{家庭や学校など、限られた範囲のネットワークを何というか。}{LAN}
\QQ{ネットワーク上で、情報やサービスを提供するコンピュータを何というか。}{サーバ}
\QQ{ネットワークどうしを中継し、データの交通整理を行う機器を何というか。}{ルータ}

\Instruction{通信のしくみに関する次の問いに答えなさい。}
\QQ{情報機器が通信するときに従う必要がある約束を何というか。}{ネットワークプロトコル}
\QQ{インターネットにつながるための標準的なネットワークプロトコルを何というか。}{TCP/IP}
\QQ{ネットワーク上で情報機器を識別するために割り当てられる数字の組み合わせを何というか。}{IPアドレス}
\QQ{数字のままではわかりにくいIPアドレスに対応させる、文字列の名前を何というか。}{ドメイン名}
\QQ{IPアドレスとドメイン名を対応させるサーバを何というか。}{DNSサーバ}

\Instruction{データ通信とサービスに関する次の問いに答えなさい。}
\QQ{インターネットで、データを効率よく送信するために分割した小さなまとまりを何というか。}{パケット}
\QQ{ネットワークを通じて、コンピュータやソフトウェアを利用できるサービスシステムを何というか。}{クラウドコンピューティング}
\QQ{情報通信ネットワークを安全に利用するための技術を何というか。}{情報セキュリティ}
\QQ{仮想空間における情報セキュリティを何というか。}{サイバーセキュリティ}
\QQ{情報の漏えいを防ぐために、データを第三者に読みにくい形に変えることを何というか。}{暗号化}

\Instruction{情報セキュリティの対策に関する次の問いに答えなさい。}
\QQ{ブラウザとサーバ間の通信の暗号化技術として使われるものを何というか。}{SSL}
\QQ{情報を盗む、機器を破壊するなど、悪意のある不正プログラムを何というか。}{マルウェア}
\QQ{コンピュータウイルスなどへの対策に利用するソフトウェアを何というか。}{セキュリティ対策ソフトウェア}
\QQ{許可された通信のみを内部に通すしくみを何というか。}{ファイアウォール}
\QQ{機器の故障などに備えて、データの予備を保存しておくことを何というか。}{バックアップ}
